\chapter*{Lời Mở Đầu}
\addcontentsline{toc}{chapter}{Lời Mở Đầu}
\markboth{Lời Mở Đầu}{Lời Mở Đầu}

\begin{quotecard}{Mở sách này để làm gì}
Cuốn này không sinh ra để mắng học sinh cho vui. Nó sinh ra để chụp đúng cái khoảnh khắc một câu chống chế nghe rất khôn, nhưng bóc thêm một lớp là thấy ngay sợ, lười, sĩ diện, hoặc mệt thật đang ngồi phía sau.

\textit{(This book is not here to scold students for fun. It is here to catch the moment when an excuse sounds clever on the surface, but underneath it hides fear, laziness, pride, or real exhaustion.)}
\end{quotecard}

\begin{truyen}{Một bản V2 rộng hơn và gai hơn}{A wider, sharper V2}
\chuhoa{B}{ản này được nới rộng theo đúng nhịp lớp học thật hơn: không chỉ có bài vở và điểm số, mà còn có làm nhóm, sống online, AI, sĩ diện trước đám đông, áp lực gia đình, chuyện chọn nghề, cả những pha drama khiến đầu óc mất sóng đúng lúc cần tỉnh nhất.}
\textit{(This edition expands into more of real school life: not just homework and grades, but group work, online behavior, AI, crowd-facing pride, family pressure, career anxiety, and the dramas that short-circuit the mind at the worst time.)}

Mỗi mục vẫn giữ cấu trúc ngắn, gọn, dễ đọc nhanh. Nhưng phần đáp trả được đẩy sắc hơn, sát tiếng nói trong lớp hơn, và cố giữ một nguyên tắc: khịa để tỉnh, không khịa cho rỗng.
\textit{(Each piece stays short and quick to read. But the replies are sharper, closer to actual classroom speech, and built on one principle: sarcasm should wake people up, not merely show off.)}

Nếu bạn là giáo viên, đây là kho tình huống để đọc ra được ẩn ý phía sau câu nói học trò. Nếu bạn là học sinh, đây là cuốn dễ làm mình phì cười trước, rồi vài phút sau mới thấy bị bóc hơi đau.
\textit{(If you are a teacher, this is a bank of situations that helps decode what lies beneath student speech. If you are a student, this book may first make you laugh, then later realize it exposed you a little too accurately.)}
\end{truyen}

\begin{baihoc}
Khịa đúng không làm người ta bé đi. Nó làm lớp ngụy biện mỏng đi để phần trưởng thành có cửa bước ra.
\textit{(The right kind of sarcasm does not shrink a person. It thins out the excuse so growth has room to appear.)}
\end{baihoc}

\begin{ghinhoanh}
\textbf{Short English to remember:}

Sharp words are useful only when they reveal truth.

If a reply only hurts and teaches nothing, it is just noise.
\end{ghinhoanh}

\begin{tuhocnhanh}
1. Em hay lý sự vì thật sự có lý, hay vì chưa dám nhận chỗ yếu? \textit{(Do you argue because you are right, or because you are not ready to face a weakness?)}

2. Một câu em từng nói để cứu sĩ diện là gì? \textit{(What is one sentence you once used just to save face?)}

3. Nếu bị bóc đúng chỗ đau, em muốn chống chế tiếp hay sửa thật? \textit{(If someone exposes a painful truth, do you want to argue further or actually fix it?)}
\end{tuhocnhanh}

\ngancach